% Generated draft: requires mathematical and source review.
\documentclass[11pt]{article}
\usepackage[margin=25mm]{geometry}
\usepackage{fontspec}
\setmainfont{Latin Modern Roman}
\usepackage{amsmath,amssymb,mathtools,unicode-math}
\setmathfont{Latin Modern Math}
\usepackage{enumitem,needspace,xurl,hyperref}
\setlength{\parindent}{0pt}
\setlength{\parskip}{5pt}
\setlength{\emergencystretch}{3em}
\newcommand{\N}{\mathbb{N}}
\newcommand{\Z}{\mathbb{Z}}
\newcommand{\Q}{\mathbb{Q}}
\newcommand{\R}{\mathbb{R}}
\newcommand{\C}{\mathbb{C}}
\newcommand{\F}{\mathbb{F}}
\newcommand{\A}{\mathbb{A}}
\newcommand{\PP}{\mathbb{P}}
\newcommand{\E}{\mathbb{E}}
\newcommand{\eps}{\varepsilon}
\newcommand{\dd}{\,\mathrm{d}}
\begin{document}
\section*{UnsolvedMath ID: 30002519}
\textbf{Source problem code:} catalogue entry
\subsection{Definitions and mathematical statement}
The supplied record uses a Bruhat-ordered indexing set, denoted here by \(W\), but does not specify its ambient group or the permitted group types. For indices \(v,w\in W\) with \(v\leq w\), let \(\mathcal R_{v,w}\) denote the open Richardson variety intended by the source. Its coordinate ring means the ring of global regular functions
\[
\C[\mathcal R_{v,w}]:=\Gamma(\mathcal R_{v,w},\mathcal O_{\mathcal R_{v,w}}).
\]
Write \(R_{v,w}\) for the cluster algebra attached to this variety in the source's intended categorical construction; this notation distinguishes the algebra from the variety \(\mathcal R_{v,w}\).

The proposed equality is
\[
\text{for every }v,w\in W\text{ with }v\leq w,\qquad
R_{v,w}=\C[\mathcal R_{v,w}].
\]
This equality is to be understood under the intended realization of the cluster algebra as regular functions on the variety. The supplied record does not specify that realization, the ambient flag variety, the precise Richardson-cell conventions, or the cluster construction and coefficient/localization conventions. Thus the display records the intended question, but these missing data prevent it from being a fully specified universal conjecture. No additional restriction on group type or choice of cluster construction is imposed here. The requested equality concerns the stated cluster algebra; an equality involving an upper cluster algebra would be a distinct assertion.
\subsection{Short English statement}
For every pair \(v\leq w\) in the intended Bruhat order, does the categorical cluster algebra \(R_{v,w}\) coincide with the ring of regular functions on the open Richardson variety \(\mathcal R_{v,w}\)? The ambient setting and cluster-algebra conventions require clarification from the source.
\subsection{Sources}
\begin{itemize}
\item Local UnsolvedMath problem registry, supplied record with permanent ID \(30002519\), title \emph{Cluster Algebras of Open Richardson Varieties}, exact source problem code \texttt{catalogue entry}; collection: Oberwolfach Reports Open Problems. The record gives the proposed year as 2014 and names no proposer.
\item Supplied report citation: \emph{Representation Theory of Quivers and Finite Dimensional Algebras} (2014), Oberwolfach Reports, pp.~453--529, DOI 10.4171/owr/2014/08. Source URL and literature-source label \emph{Oberwolfach Report 08/2014}: \url{https://doi.org/10.4171/owr/2014/08}. The supplied background describes these reports as records of mathematical workshops and their open-problem discussions.
\item Supplied literature-source label \emph{Leclerc conjecture in type A}: \url{https://arxiv.org/abs/2210.13302}.
\item Supplied literature-source label \emph{Open Richardson upper cluster algebras}: \url{https://arxiv.org/abs/2207.10184}.
\item Literature assessment dated exactly \texttt{2026-08-21}: the supplied curation labels the problem partially solved and states that the equality between the categorical cluster algebra and the coordinate ring is proved in type A, that upper-cluster-algebra results and special cases are known more generally, and that the full stated generality remains open. This is supplied dated curation derived from the OWR material, not an independently verified literature assessment; the time-sensitive status and cited claims require independent verification.
\end{itemize}
\end{document}
